\documentclass[11pt, a4paper]{article}

\usepackage[a4paper, top=2.5cm, bottom=2.5cm, left=2cm, right=2cm]{geometry}
\usepackage{amsmath}
\usepackage{amssymb}
\usepackage[colorlinks=true, linkcolor=blue, urlcolor=blue, citecolor=blue]{hyperref}

\title{Theory of Everything - Volume 49 \\ \Large Quantum Reference Frames}
\author{Omega Protocol}
\date{\today}

\begin{document}

\maketitle

\section*{Layman's Summary}
There is no "God's eye view" of the universe. Volume 49 proves that reality is entirely relative. Your perspective of the universe depends fundamentally on how you are entangled with it. We resolve deep quantum paradoxes by proving that objective reality is just the overlapping consensus of all subjective viewpoints.

\section{Relational Quantum Mechanics (Axiom 52)}
We establish that quantum states do not exist in a vacuum. A particle doesn't simply "have" a location or a speed; it only has a location *relative to* a specific observer node. The entire structure of the universe is built on these relative informational relationships.

\section{Wigner's Friend Resolved (Theorem)}
Imagine a scientist in a box making a measurement, while you stand outside the box. To the scientist, the measurement is done. To you, the scientist is still in a fuzzy quantum superposition. Who is right? We prove that both are right. Contradictions only happen if you assume a universal absolute truth. When you account for the network geometry, the paradox vanishes.

\section{Gauge Invariance of Reality (Theorem)}
If everyone sees a slightly different quantum reality, does objective reality even exist? We prove that it does. The objective, physical universe is the "gauge-invariant intersection"—the core mathematical structure that remains perfectly consistent across all subjective reference frames. Reality is the mathematical overlap of everyone's perspective.

\end{document}
